\documentclass[a4paper,12pt]{article}

\usepackage[T2A]{fontenc}
\usepackage[utf8]{inputenc}
\usepackage[russian]{babel}
\usepackage{amsmath,amssymb,mathtools}
\renewcommand{\familydefault}{\sfdefault}
\usepackage{icomma}
\usepackage[a4paper,margin=20mm,headheight=25pt,headsep=8mm,footskip=12mm]{geometry}
\usepackage{microtype}
\usepackage{tikz}
\usetikzlibrary{calc,arrows.meta,positioning,shapes.geometric}
\usepackage{pgfplots}
\pgfplotsset{compat=1.18}
\usepackage{tabularx,booktabs}
\usepackage{enumitem}
\usepackage{parskip}
\usepackage{fancyhdr}
\usepackage{setspace}
\usepackage{xcolor}

\definecolor{accent}{HTML}{008080}
\definecolor{darkaccent}{HTML}{004D4D}
\definecolor{textgray}{HTML}{333333}
\definecolor{softgray}{HTML}{6B7280}
\definecolor{lightaccent}{HTML}{D9ECEC}

\onehalfspacing
\setlength{\parindent}{0pt}
\setlength{\emergencystretch}{2em}
\setlist[itemize]{leftmargin=6mm,itemsep=0.3em,topsep=0.35em}
\setlist[enumerate]{leftmargin=7mm,itemsep=0.35em,topsep=0.35em}

\pagestyle{fancy}
\fancyhf{}
\fancyhead[L]{\small\textcolor{softgray}{Физика, 7 класс}}
\fancyhead[R]{\small\textcolor{softgray}{Кирилл Зиновьев}}
\fancyfoot[C]{\small\textcolor{softgray}{\thepage}}
\renewcommand{\headrulewidth}{0.3pt}
\renewcommand{\headrule}{%
  \hbox to\headwidth{%
    \color{lightaccent}\leaders\hrule height \headrulewidth\hfill
  }%
}

\newcommand{\sectionline}{%
  \par\vspace{0.15em}%
  \noindent\textcolor{lightaccent}{\rule{\linewidth}{0.7pt}}%
  \par\vspace{0.35em}%
}

\newcommand{\key}[1]{\textcolor{darkaccent}{\textbf{#1}}}
\newcommand{\unit}[1]{\,\text{#1}}

\tikzset{
  startstop/.style={
    ellipse,
    draw=darkaccent,
    line width=0.9pt,
    align=center,
    minimum width=35mm,
    minimum height=10mm,
    inner sep=4pt,
    font=\small
  },
  process/.style={
    rounded corners=3pt,
    rectangle,
    draw=accent,
    line width=0.9pt,
    align=center,
    text width=78mm,
    minimum height=12mm,
    inner sep=5pt,
    font=\small
  },
  decision/.style={
    diamond,
    aspect=2.2,
    draw=darkaccent,
    line width=0.9pt,
    align=center,
    text width=42mm,
    inner sep=2.5pt,
    font=\small
  },
  flowarrow/.style={
    -{Latex[length=3mm,width=2mm]},
    line width=0.9pt,
    draw=textgray
  },
  branch/.style={
    font=\small,
    fill=white,
    inner sep=1pt,
    text=textgray
  }
}

\begin{document}

\thispagestyle{empty}

\begin{center}
\vspace*{2.2cm}

{\Large\textcolor{darkaccent}{\textbf{Чек-лист}}}\par
\vspace{0.5cm}

{\LARGE\textbf{Равномерное движение: единицы, размерности и графики}}\par
\vspace{1.0cm}

{\large Кирилл Зиновьев, 7 класс}\par
\vspace{0.35cm}

{\normalsize Физика \quad\textbullet\quad Тренировочный блок: Путь А}\par

\vfill

{\large \today}\par
\vspace{0.9cm}

{\normalsize Лёвин Артём Александрович}\par
{\normalsize эксклюзивно для Кирилла Зиновьева}\par

\vspace{1.5cm}

\begin{tikzpicture}[remember picture,overlay]
  \draw[
    accent!55,
    line width=1.1pt
  ]
  ([xshift=22mm,yshift=18mm]current page.south west)
  .. controls
  ([xshift=62mm,yshift=30mm]current page.south west)
  and
  ([xshift=-62mm,yshift=10mm]current page.south east)
  ..
  ([xshift=-22mm,yshift=22mm]current page.south east);
\end{tikzpicture}
\end{center}

\clearpage

\section*{1. Основа темы}
\sectionline

\begin{itemize}

  \item
  \key{Равномерное движение} --- движение с постоянной по модулю скоростью.
  За равные промежутки времени тело проходит равные пути.

  \item
  Для движения из начальной точки отсчёта пути:
  \[
    s=vt,
    \qquad
    v=\frac{s}{t}\quad (t>0),
    \qquad
    t=\frac{s}{v}\quad (v>0).
  \]

  \item
  В задачах на путь считаем
  \[
    s\ge0,\qquad t\ge0,\qquad v\ge0.
  \]

  \item
  \key{По умолчанию удобно приводить данные к СИ:}
  путь --- в метры,
  время --- в секунды,
  скорость --- в метры в секунду.

  \item
  Иногда можно оставить другие \emph{согласованные} единицы,
  если это упрощает расчёт.
  Например,
  $6\unit{м}$ и $0{,}3\unit{м/год}$
  удобно использовать напрямую,
  если нужен ответ в годах.

\end{itemize}

\section*{2. Перевод единиц}
\sectionline

\renewcommand{\arraystretch}{1.35}

\begin{table}[h]
\centering
\begin{tabularx}{\linewidth}{@{}X X@{}}
\toprule
\textbf{Переход} & \textbf{Правило} \\
\midrule
минуты $\to$ секунды
&
умножить на $60$
\\
часы $\to$ секунды
&
умножить на $3600$
\\
километры $\to$ метры
&
умножить на $1000$
\\
сантиметры $\to$ метры
&
разделить на $100$
\\
миллиметры $\to$ метры
&
разделить на $1000$
\\
$\text{км/ч}\to\text{м/с}$
&
разделить числовое значение на $3{,}6$
\\
\bottomrule
\end{tabularx}
\end{table}

Почему для скорости появляется число $3{,}6$:

\[
  1\unit{км/ч}
  =
  \frac{1000\unit{м}}{3600\unit{с}}
  =
  \frac{1}{3{,}6}\unit{м/с}.
\]

Поэтому, например,

\[
  72\unit{км/ч}
  =
  \frac{72}{3{,}6}
  =
  20\unit{м/с}.
\]

\section*{3. Быстрые примеры}
\sectionline

\[
  3\unit{мин}
  =
  3\cdot60
  =
  180\unit{с},
  \qquad
  2\unit{ч}
  =
  2\cdot3600
  =
  7200\unit{с}.
\]

\[
  3{,}5\unit{км}
  =
  3500\unit{м},
  \qquad
  3{,}5\unit{см}
  =
  0{,}035\unit{м},
  \qquad
  3{,}5\unit{мм}
  =
  0{,}0035\unit{м}.
\]

\[
  1{,}4\unit{мм/с}
  =
  \frac{1{,}4}{1000}\unit{м/с}
  =
  0{,}0014\unit{м/с}.
\]

\section*{4. Размерность как способ самопроверки}
\sectionline

Размерности можно преобразовывать вместе с числами. Например,

\[
  \frac{\text{м}}{\text{с}}\cdot\text{с}
  =
  \text{м}.
\]

Если после вычисления пути в размерности остались секунды,
значит в записи есть ошибка.

\textbf{Пример.}
Улитка движется со скоростью
$1{,}4\unit{мм/с}$
в течение
$5\unit{мин}$.
Найдём путь.

Сначала переведём время:

\[
  5\unit{мин}
  =
  300\unit{с}.
\]

\textbf{Способ 1: привести скорость к СИ.}

\[
  v
  =
  0{,}0014\unit{м/с},
  \qquad
  s
  =
  vt
  =
  0{,}0014\cdot300
  =
  0{,}42\unit{м}.
\]

\textbf{Способ 2: сначала получить путь в миллиметрах.}

\[
  s
  =
  1{,}4\unit{мм/с}\cdot300\unit{с}
  =
  420\unit{мм}
  =
  0{,}42\unit{м}.
\]

Оба способа корректны.
Важно понимать,
в какой единице получен промежуточный результат.

\section*{5. Поиск времени}
\sectionline

Средняя скорость роста дерева равна
$0{,}3\unit{м/год}$,
а прирост высоты --- $6\unit{м}$.

Тогда

\[
  t
  =
  \frac{s}{v}
  =
  \frac{6\unit{м}}{0{,}3\unit{м/год}}
  =
  20\text{ лет}.
\]

Проверка единицы:
\[
  \frac{\text{м}}{\text{м}/\text{год}}
  =
  \text{год}.
\]

\clearpage
\thispagestyle{empty}

\begin{center}

{\Large
\textcolor{darkaccent}{
\textbf{Алгоритм: расчёт при равномерном движении}
}}

\vspace{8mm}

\begin{tikzpicture}[node distance=7mm and 22mm]

\node[startstop] (start)
{Начало};

\node[
  process,
  below=of start
] (data)
{Выпишите известные величины, искомую величину и единицы};

\node[
  process,
  below=of data
] (formula)
{Выберите формулу: $s=vt$, $v=s/t$ или $t=s/v$};

\node[
  decision,
  below=9mm of formula
] (units)
{Единицы согласованы и приведут к нужной единице ответа?};

\node[
  process,
  below=13mm of units
] (substitute)
{Подставьте значения вместе с единицами};

\node[
  process,
  left=24mm of substitute,
  text width=58mm
] (convert)
{Переведите необходимые величины в согласованные единицы};

\node[
  process,
  below=of substitute
] (calculate)
{Выполните вычисление и сократите размерности};

\node[
  decision,
  below=9mm of calculate
] (check)
{Получена требуемая единица?};

\node[
  process,
  below=13mm of check
] (answer)
{Запишите ответ с единицей измерения};

\node[
  process,
  left=24mm of answer,
  text width=58mm
] (convertanswer)
{Переведите результат в требуемую единицу};

\node[
  startstop,
  below=of answer
] (end)
{Конец};

\draw[flowarrow]
(start)
--
(data);

\draw[flowarrow]
(data)
--
(formula);

\draw[flowarrow]
(formula)
--
(units);

\draw[flowarrow]
(units)
--
node[branch,right]{да}
(substitute);

\draw[flowarrow]
(units.west)
-|
node[branch,pos=0.22]{нет}
(convert.north);

\draw[flowarrow]
(convert.east)
--
(substitute.west);

\draw[flowarrow]
(substitute)
--
(calculate);

\draw[flowarrow]
(calculate)
--
(check);

\draw[flowarrow]
(check)
--
node[branch,right]{да}
(answer);

\draw[flowarrow]
(check.west)
-|
node[branch,pos=0.22]{нет}
(convertanswer.north);

\draw[flowarrow]
(convertanswer.east)
--
(answer.west);

\draw[flowarrow]
(answer)
--
(end);

\end{tikzpicture}

\end{center}

\clearpage

\section*{6. График зависимости пути от времени}
\sectionline

При постоянной скорости

\[
  s(t)=vt.
\]

\begin{itemize}

  \item
  $t$ --- \key{независимая переменная}:
  её значения выбирают сами;
  она откладывается по горизонтальной оси абсцисс.

  \item
  $s$ --- \key{зависимая переменная}:
  её вычисляют по выбранному $t$;
  она откладывается по вертикальной оси ординат.

  \item
  $v$ --- постоянная величина для данного движения.

\end{itemize}

Для

\[
  v_1=7\unit{м/с},
  \qquad
  v_2=3\unit{м/с}
\]

получаем

\[
  s_1=7t,
  \qquad
  s_2=3t.
\]

Для каждой прямой достаточно двух различных точек.
Удобно взять
$t=0$
и
$t=1\unit{с}$:

\[
  (0;0),\ (1;7)
  \qquad\text{и}\qquad
  (0;0),\ (1;3).
\]

\begin{center}

\begin{tikzpicture}

\begin{axis}[
  width=0.9\linewidth,
  height=7.8cm,
  axis lines=middle,
  xmin=0,
  xmax=5.2,
  ymin=0,
  ymax=37,
  xlabel={$t,\ \text{с}$},
  ylabel={$s,\ \text{м}$},
  xtick={0,1,2,3,4,5},
  ytick={0,5,10,15,20,25,30,35},
  grid=both,
  minor tick num=0,
  clip=false,
  samples=200,
  legend style={
    at={(0.97,0.05)},
    anchor=south east,
    draw=none,
    fill=none,
    font=\small
  }
]

\addplot[
  accent,
  line width=1.2pt,
  domain=0:5
]
{7*x};

\addlegendentry{$s_1=7t$}

\addplot[
  darkaccent!75,
  line width=1.2pt,
  dashed,
  domain=0:5
]
{3*x};

\addlegendentry{$s_2=3t$}

\end{axis}

\end{tikzpicture}

\end{center}

\textbf{Как читать график:}

\begin{itemize}

  \item
  при $t=0$ путь равен нулю,
  поэтому обе прямые проходят через начало координат;

  \item
  чем больше скорость,
  тем круче поднимается график $s(t)$;

  \item
  чтобы найти путь в выбранный момент времени,
  берут значение $t$ на горизонтальной оси
  и читают соответствующее $s$ по вертикальной.

\end{itemize}

\clearpage
\thispagestyle{empty}

\begin{center}

{\Large
\textcolor{darkaccent}{
\textbf{Алгоритм: построение графика $s(t)$}
}}

\vspace{8mm}

\begin{tikzpicture}[node distance=7mm and 22mm]

\node[startstop] (startg)
{Начало};

\node[
  process,
  below=of startg
] (writeg)
{Запишите зависимость $s=vt$};

\node[
  process,
  below=of writeg
] (axesg)
{Проведите оси и подпишите: по горизонтали $t$, по вертикали $s$};

\node[
  process,
  below=of axesg
] (timesg)
{Выберите не менее двух удобных значений времени};

\node[
  process,
  below=of timesg
] (calcg)
{Для каждого $t$ вычислите $s$ по формуле};

\node[
  process,
  below=of calcg
] (pointsg)
{Нанесите точки $(t;s)$ на координатную плоскость};

\node[
  process,
  below=of pointsg
] (lineg)
{Проведите прямую через построенные точки};

\node[
  decision,
  below=9mm of lineg
] (checkg)
{Оси, единицы и масштаб подписаны корректно?};

\node[
  startstop,
  below=13mm of checkg
] (endg)
{График готов};

\node[
  process,
  left=24mm of endg,
  text width=58mm
] (fixg)
{Исправьте подписи осей, единицы или масштаб};

\draw[flowarrow]
(startg)
--
(writeg);

\draw[flowarrow]
(writeg)
--
(axesg);

\draw[flowarrow]
(axesg)
--
(timesg);

\draw[flowarrow]
(timesg)
--
(calcg);

\draw[flowarrow]
(calcg)
--
(pointsg);

\draw[flowarrow]
(pointsg)
--
(lineg);

\draw[flowarrow]
(lineg)
--
(checkg);

\draw[flowarrow]
(checkg)
--
node[branch,right]{да}
(endg);

\draw[flowarrow]
(checkg.west)
-|
node[branch,pos=0.22]{нет}
(fixg.north);

\draw[flowarrow]
(fixg.west)
--
++(-12mm,0)
|-
(axesg.west);

\end{tikzpicture}

\end{center}

\clearpage

\section*{7. Типичные ошибки}
\sectionline

\begin{enumerate}

  \item
  Смешать минуты и секунды,
  километры и метры
  внутри одного вычисления
  без явного перевода.

  \item
  Перепутать переходы:
  $\text{см}\to\text{м}$ --- деление на $100$,
  $\text{мм}\to\text{м}$ --- деление на $1000$.

  \item
  Забыть единицу у ответа.

  \item
  Получить путь,
  но оставить в размерности лишнюю секунду:
  это сигнал ошибки в формуле или единицах.

  \item
  Перепутать оси графика $s(t)$:
  время $t$ идёт по горизонтали,
  путь $s$ --- по вертикали.

  \item
  Выбирать значения $s$ произвольно.
  Свободно выбирают $t$,
  а $s$ затем вычисляют по формуле.

  \item
  При использовании
  $v=s/t$
  забыть,
  что $t>0$;
  при использовании
  $t=s/v$ ---
  что $v>0$.

\end{enumerate}

\section*{8. Самопроверка перед ответом}
\sectionline

\begin{itemize}[label={$\square$}]

  \item
  Выписаны известные и искомая величины.

  \item
  Единицы согласованы между собой.

  \item
  Выбрана формула для нужной величины.

  \item
  Размерности приводят к единице искомой величины.

  \item
  Численный ответ разумен по порядку величины.

  \item
  Ответ записан с единицей.

  \item
  Для графика подписаны обе оси и единицы измерения.

\end{itemize}

\clearpage

\section*{Тренировочный блок --- Путь А}
\sectionline

\textbf{Задание 1. Перевод единиц.}

Переведите в СИ:

\begin{enumerate}[label=\alph*)]

  \item
  $4{,}5\unit{мин}$;

  \item
  $1{,}25\unit{ч}$;

  \item
  $2{,}8\unit{км}$;

  \item
  $6{,}4\unit{см}$;

  \item
  $9\unit{мм}$;

  \item
  $54\unit{км/ч}$.

\end{enumerate}

\textbf{Задание 2. Путь.}

Муравей движется со скоростью
$2{,}5\unit{мм/с}$.
Какой путь в метрах он пройдёт за
$8\unit{мин}$?

\textbf{Задание 3. Время.}

Средняя скорость роста растения равна
$0{,}24\unit{м/год}$.
За сколько лет его высота увеличится на
$3{,}6\unit{м}$?

\textbf{Задание 4. Согласование единиц.}

Тело движется равномерно со скоростью
$18\unit{км/ч}$
в течение
$40\unit{с}$.
Найдите путь в метрах.

\textbf{Задание 5. График движения.}

Два тела одновременно начинают движение из одной точки.
Их скорости постоянны:

\[
  v_1=4\unit{м/с},
  \qquad
  v_2=6\unit{м/с}.
\]

\begin{enumerate}[label=\alph*)]

  \item
  Запишите зависимости
  $s_1(t)$
  и
  $s_2(t)$.

  \item
  Постройте оба графика
  на одной координатной плоскости
  для
  \[
    0\le t\le5\unit{с}.
  \]

  \item
  Определите путь каждого тела
  через
  $3\unit{с}$.

  \item
  Укажите,
  какой график круче,
  и объясните это через значение скорости.

\end{enumerate}

\clearpage

\section*{Ответы к тренировочному блоку}
\sectionline

\renewcommand{\arraystretch}{1.45}

\begin{table}[h]
\centering
\begin{tabularx}{\linewidth}{@{}p{0.08\linewidth}X@{}}

\toprule
\textbf{№} & \textbf{Ответ} \\
\midrule

1
&
а) $270\unit{с}$;
б) $4500\unit{с}$;
в) $2800\unit{м}$;
г) $0{,}064\unit{м}$;
д) $0{,}009\unit{м}$;
е) $15\unit{м/с}$.
\\

2
&
$1{,}2\unit{м}$.
\\

3
&
$15$ лет.
\\

4
&
$200\unit{м}$.
\\

5
&
а) $s_1=4t$, $s_2=6t$;
б) прямые через начало координат;
для проверки можно использовать точки
$(5;20)$
и
$(5;30)$
соответственно;
в) $12\unit{м}$ и $18\unit{м}$;
г) круче график $s_2=6t$,
поскольку
$6\unit{м/с}>4\unit{м/с}$.
\\

\bottomrule

\end{tabularx}
\end{table}

% :chatgpt-content-reference{index="0"}

\end{document}